\section{PLANTEAMIENTO DEL PROBLEMA}

\subsection{PROBLEMA CENTRAL}
Los nuevos algoritmos criptogr\'{a}ficos resistentes a ataques cu\'{a}nticos generan datos de identificaci\'{o}n y firma considerablemente m\'{a}s grandes que los esquemas cl\'{a}sicos, lo que puede afectar la velocidad y eficiencia de los sistemas que los utilizan.

¿De qué manera el uso de criptografía post-cuántica basada en retículos afecta el rendimiento computacional y la viabilidad de los sistemas de votación electrónica cuando operan bajo una alta concurrencia de usuarios?


\subsection{PROBLEMAS SECUNDARIOS}

\begin{itemize}
    \item El aumento en la latencia total y retrasos al procesar el sufragio, provocado por el volumen de datos de las firmas FIPS~204 y la complejidad computacional de las restricciones en los circuitos de conocimiento cero (ZKP). (Lopez et al, 2025; Chaurasia, Sharma y Sharma, 2025)

    \item El incremento en los costos transaccionales (unidades de gas) dentro de la M\'{a}quina Virtual de Ethereum (EVM), debido al esfuerzo de c\'{o}mputo necesario para procesar las operaciones de la criptograf\'{i}a basada en ret\'{i}culos. (Ezeh, Oparah, Gado, Gbaraba y Adeleke, 2025; Finogina y Herranz, 2023)

    \item La saturaci\'{o}n y p\'{e}rdida de rendimiento f\'{i}sico (CPU y memoria RAM) en los nodos o servidores convencionales al gestionar m\'{u}ltiples solicitudes concurrentes de votantes bajo condiciones de alta demanda. (FS-ISAC, 2023; Lopez et al, 2025)
\end{itemize}
